% New independent scoped review. No frozen author artifact is edited.
\subsection{Independent review of the affine-parent fourth lower integral}
\label{sec:independent-fourth-chord-review}

The reviewed object is the ideal fourth integral $I_4$ for the same
causal $C_3=\Phi_3$, followed by the all-licensed-mask maximum $C_4$.
It is not an increment to a previous numerical estimate of $I_4$.
The reviewed sources are \texttt{lower-start-fourth-affine-chord-proof.tex}
and \texttt{certify\_lower\_start\_fourth\_chord.py}.

\paragraph{Parent symmetry and history.}
On the retained ordered domain $t_1<w=19/125<25/128$, the redundant
singleton $v(t_3)$ is bounded by $v(t_1)$, since $v$ is increasing
on this interval. The maximum of the three $W$ terms is precisely
$\Lambda-3\min_jt_j/8$. Thus the seven-subset symmetric expression
is the same function as the original ordered eight-branch policy.
The middle affine formula for $v$ is redundant: at the crossing
$25/128$ of the outer two formulas it lies strictly above them,
and its slope lies between their slopes. This argument holds on
both sides of that crossing. The symmetric extension does not
assert that a gate invalid on the ordered domain becomes valid.
Admission still uses the original old history.

For each retained box let $L_3\le C_3\le U_3$ be the exact range
bounds. The inequalities $t_{1,+}+3t_{2,+}<\nu$ and
$r<(L_3-\sum t_{j,+})/3$ imply every needed old history,
including old residual argument above two, for every ordering
permitted by equal outer intervals. A separate affine majorant is
used only to bound the subtracted density. Replacing $L_3$ in the
history test by that majorant would be invalid; the code does not
make this replacement.

\paragraph{Branch majorant and all remaining orders.}
For every min-of-affine parent branch the shift construction
exhibits at least one affine member bounded above by the chosen
majorant throughout the box. Taking the maximum over branches
therefore preserves the inequality $C_3\le U$. The independent
replay rechecks this branch by branch, rather than trusting a
zero shift or the branch selected at the approximate mean.

For a candidate fourth subset, each possible ordering of its
remaining old slots is enumerated. Every required square gate is
an affine expression consisting of the selected old sum, $r$,
the previous remaining slots and twice the current remaining slot.
Both affine pieces of $v$ dominate this expression on the retained
whole box and last-variable interval. The further test
$V-S_3\ge3r$ keeps the new residual argument in $[2,\infty)$.
Equality at a last-variable endpoint has zero measure; it is
handled by the existing buffered finite approximation, not claimed
as a strict finite-$N$ admission. No gate is checked only at an
approximate center.

\paragraph{Shape, integration and signs.}
Ordinary $\bar f$ is increasing and concave on $[2,\infty)$.
Consequently $-\bar f((U-S_3-r)/r)$ is convex in the old
coordinates, and Jensen gives a lower bound for the subtracted
old contribution at its logarithmic mean. Since $V$ is a minimum
of affine functions, $\bar f((V-S_3-r)/r)$ is concave on the
entire retained box. Iterated one-coordinate chord inequalities
give its eight-corner multilinear lower bound, including across
the tent's kink. This uses neither an extrapolation below two
nor a bound on higher derivatives.

The primitive formula is
\[
 \int_l^u\bar f((D+Br-r)/r)\,\frac{dr}{r^2}
 =\frac{H(D/l+B)-H(D/u+B)}D,
 \qquad H'(s)=\bar f(s-1),\quad D>0.
\]
All primitive arguments are checked in $[3,32]$. The old mean
uses an upper enclosure of $E(U-S_3)$ and hence an upper old
integral; each positive new corner integral uses a lower
nonnegative weight. Independent lower-minus-upper arithmetic
then has the claimed direction. Taking a positive part is
permitted only because the actual ideal fourth gap is
nonnegative. It does not assert pointwise positivity of the
signed estimator.

\paragraph{Disjointness and multiplicities.}
For repeated outer intervals the actual ideal integrand is
symmetric. Its ordered-sector integral equals the full-box
integral divided by the product of the repetition factorials.
The numerical estimator may be asymmetric, provided it bounds
the actual integrand throughout the full box; the all-order
gate tests give precisely that property. Splitting a group of
$m$ equal intervals into its $m+1$ possible ordered high/low
multiplicities gives a disjoint sector partition. No binomial
factor is to be added to the factorial sector weight.

The independent replay reconstructs the complete finite history
partition and checks uniqueness of every saved positive parent.
The last-order strip has $r\ge a_3$, whereas the main slices
have $r<a_3$. Its mass formula follows by integrating
$\log(b_3/r)/r^2$. The common-corner gap is used there only
with fixed old and new caps. The two main parent estimators
are maximized, and the disjoint strip is then added once.
No two candidate masks are summed over the same region.

\paragraph{Numerical scope.}
The new replay independently reconstructs all 37291 history
leaves and all 10384 saved positive parent bounds. It imports
neither the author's chord source nor its Jensen and slice
helpers. It deliberately shares the previously independently
audited ordinary mesh-65536, $112$-bit delay table. Its interval
and logarithm implementation is separate. Thus it is an
independent integration and geometry replay with an explicitly
shared ordinary backend, not a second independent construction
of that backend. Its JSON reports the exact total, any individual
floor comparisons that are weaker because of outward rounding,
and the actual comparison scope. A valid slightly smaller total
can be used directly instead of treating the author's exact
integer as independently covered.

No analytic contradiction was found in the reviewed construction.
The passage from this ideal integral to one common finite causal
policy, and then to arithmetic counts for all sufficiently large
even $N$, still uses the separately reviewed finite approximation
and distribution statements. Those statements are not consequences
of this numerical replay.
